\chapter[Sermon 62. Christ Stands upon Mount Zion, Having His Church: and Is Described by Notes, Which and What Shall Be the Sheep of Christ.]{Christ Stands upon Mount Zion, Having His Church: and Is Described by Notes, Which and What Shall Be the Sheep of Christ.}
\label{sermon:062}
\markright{Sermon 62. Christ Stands upon Mount Zion, Having His Church: and ...}

And I looked, and behold, a Lamb stood on Mount Zion, with him one hundred forty-four thousand, having the name of his Father written on their foreheads. And I heard a voice from heaven, like the sound of many waters, and like the voice of a great thunder. And the voice that I heard was like the sound of harpers playing on their harps. And they sang as it were a new song before the throne, and before the four living creatures, and the elders, and no one could learn that song except the one hundred forty-four thousand who were redeemed from the earth. These are they who were not defiled with women, for they are virgins. These follow the Lamb wherever he goes. These were redeemed from among people, the first fruits to God and to the Lamb, and no deceit was found in their mouths. For they are without blemish before the throne of God.

\index{Antichrist}\index{Rome}Like as he has hitherto mixed joyful things with sorrowful, and annexed a consolation to most hard and cruel chances: so now he adioyneth also to the tyranny of the Roman Empire an exposition having both a consolation and an exhortation most grave and weighty. Undoubtedly, by the description of the Romish tyranny and reign of Antichrist, it might have seemed that the Church and the preaching of the Gospel had been utterly lost, and that ungodliness should have triumphed forever; he declares therefore by a most excellent vision how Christ shall reign notwithstanding in his chosen, and shall overcome, and shall have his church continually, and that right famous. He describes what the elect shall be. He adds that the preaching of the Gospel cannot be so oppressed, but that it shall rather be preached with great constancy throughout all the world. And that Rome also shall fall, and all the ungodly be punished. He exhorts therefore most earnestly that we have nothing to do with Antichrist, lest also we be made partakers of his damnation. And to the intent there might want nothing that concerned a full comfort, he adds that thing which may chiefly confirm the minds of all the godly even in the greatest dangers, how they that die in Christ do flit straightways from corporal death unto life everlasting. Which finished, he turns to the description of the punishment to be taken assuredly of the Antichristians. Wherefore if the Books of the Gospel and New Testament are to be esteemed for the manifold description of Christ, and of salvation by him obtained for the faithful, if they are to be esteemed of the comfort and preaching of the gospel: this is doubtless a book most gospel-like, as that which by a continual tenor to perilous things annexeth consolation.

John therefore sees the Lamb standing upon Mount Zion. Christ therefore sleeps not; he is not ignorant of the perils and conflicts of his church, but he stands prepared to aid and succor his people. He stands as an invincible king, whom neither the Dragon, nor the old beast, nor the new beast has overthrown. For I have told you often, especially in chapter 5, that by the Lamb is understood Christ. For he is the Lamb and price of our redemption until the judgment, but laying aside the office of an intercessor, he shall be a most severe and also a most holy judge. And Christ stands, not in the sand, as did the Dragon, but on a Mount, and that upon Mount Zion. Mount Zion was a figure of Christ's kingdom, as appears plainly in Psalm 2 and Isaiah 2. And the kingdom of Christ is the church, both triumphant and militant; therefore, in the fellowship of saints, stands Christ, the joy and glory of those in heaven, and the life and helper of those who fight yet on Earth. Let us believe therefore that in the Antichristian persecutions Christ will never fail his faithful, as he is known never to have failed the old saints under the old Roman Empire afflicted. For this consolation serves chiefly for us, who are vexed by Antichrist, and served for them also, who suffered martyrdom under the old Roman Empire. Neither is there any doubt but that they confirmed themselves with this in the greatest persecutions.

\index{Papacy}It is most comforting to know that the Lamb is not alone, but has with him one hundred forty-four thousand—a vast and complete church. Therefore, even though the beast rages and slays the confessors of Christ, there will always be a church that cannot be uprooted from the earth. He sets a specific number for something that appears uncertain, yet is certain and determinate: for the number of those who will be saved will seem small compared to those who worship the beast and perish. Nevertheless, we understand that the number of those who will comprise the body of the Church, under their head Christ, will be greatest, even at the time when the Pope and all the agents of Antichrist have unleashed all their fury. I have spoken of this number of the elect in Chapter 7, where the same number is given.

\index{John, Saint}\index{Turks}And just as those who oppose Christ bear the mark of Antichrist on the right hand and foreheads, so truly the sheep of Christ, which will be the church, the bride of Christ, under their head Christ, will have their mark also on their foreheads, namely, the name of the Father of the Lamb. For “Eius” is to be referred to the Lamb. And he speaks not of an external mark, which should be printed on their foreheads, but of the mark of their minds. This is faith, the sign of all God’s children. And faith in the Father and the Son, which is not without the Holy Spirit. And how could you believe that almighty God is your father, unless you understand the same to be obtained through the Son? This faith, therefore, is here understood to be a Christian faith, not a Jewish or Turkish faith, which yet confesses God to be the Father. But since they do not have the Son, as Saint John says in his epistle, they neither have the Father. Therefore, the true members of the church of Christ, the true sheep, do believe that they have a merciful Father through the Son, by whom they know that the Father, being pacified, has given all things of life and of salvation in his Son. Those who do not seek salvation and all goodness in the only mediator, the Son of God, do not, without doubt, have the true mark of the children of God on their foreheads. At this day all will claim to be Christian, but neglecting Christ, they depend wholly on saints. Therefore, their faith is not the true mark of the children of God. No, they neither know the Father nor the Son. And therefore they persecute those who cleave wholly to the Father through the Son. And saying that Christ is with his church, what need has the church of a vicar? Certainly, it cannot be the true church, which has a vicar of Christ, for then it lacks Christ, whom the true church cannot lack.

It was not enough for the Apostle to have simply said that the church is united with Christ; unless he had also added, with many words, how he has seen the church afflicted and how she has behaved when the beasts torment her, then truly, so that we may learn from it, what the hope of the saints is in the greatest dangers, and of what sort we ought to be in persecutions and temptations.

First he hears a voice from heaven, as the voice of many waters. Waters in the Scriptures many times signify people. We understand therefore hereby that the church shall be populous and speaking, intending to dissemble nothing, but freely to profess Christ. And therefore he hears also the sound of a great thunder. For the church receives from heaven power to preach and show forth the Gospel gravely, though the world’s bowels burst. And verily, concerning the fervent and constant preaching of the gospel, John and James are called with Mark, sons of thunder. And concerning the preaching of the gospel, more shall follow afterward. He hears moreover a melodious harmony of men singing to their harps, and singing as it were a new song. This is chiefly referred to the saints in heaven, singing eternal praises to God; secondly, to the saints living here yet in earth, who also offer unto God continually praises and thanksgivings. Therefore, however their hearts be made sorrowful in perils and adversities, yet their spirit rejoices in the Lord. For no man could learn that same song, except the elect. For just as none of the heavenly dwellers can express or understand the excellency of the joys of the life to come and the praises of God, except he dwell among the heavenly inhabitants and be partakers of the most godly life; so except any man living yet here in earth be regenerated, he neither sees how great is the felicity of the faithful, nor can he justly esteem the praises which they offer unto God. Touching the new song, I have spoken in chapter 5. And certainly, to worldly men, the things seem as they were new, which the faithful bring forth of God’s word.

Now John describes what kind of people the sheep of Christ will be, those who will remain in the church of Christ, disregarding the rage of the beasts. To them also the mark of their fathers’ name in their foreheads is explained.

We shall also perceive what the true marks of the faithful are. First, they are redeemed from the earth. Doubtless, all of us bearing the earthly image of the earthly man were sold into sin, for which cause we are also subject to condemnation. But the Lord has bought us with the price of redemption, paid upon the cross, so that now we are reshaped into the image of the heavenly man, and are adopted as the children of God. Concerning this redemption, the Apostle has spoken in 1 Corinthians 7 and to the Romans 3, and in other places. Peter also, in 1 Peter 1. And because the faithful know themselves to be bought and adopted by Christ for the heavenly inheritance, they are eager to serve their Redeemer only, and inseparable cleave to him.

\index{Hebrews, Epistle to the}Moreover, they are virgins, not defiled by association with women. Some who discuss these matters torment themselves, fearing that something here might seem to diminish the value of holy matrimony, by which, without doubt, the Apostle testifies in 1 Corinthians 7 and to the Hebrews—no one is defiled. I am ashamed to bring forth the trivialities of the Papists. For who could hear the most corrupt reasoning about purity? They will use this to justify and defend their single lives; but all men see, except they be more blind than moles, what filthiness has been committed and is committed daily under the guise of this ungracious and most unclean singleness. But the Lord speaks nothing at this present of physical marriage, but rather spiritual. For it is evident that the apostles, as the bride leaders of our Savior, have brought the church to our Savior as a chaste virgin, which has not had dealings with any strange or foreign woman—that is to say, which is not defiled by the participation of evil doctrine. Read Solomon reasoning about that woman seriously in chapter 4 of Proverbs. Read moreover the Apostle in 2 Corinthians 11, teaching exceedingly well that the faithful are an undefiled virgin, the spouse of Christ. The faithful therefore who lived under the tyranny of the beasts received no strange doctrine of idols and of other profane cults, neither do they at this day admit the Popish infection, but keep their maidenly minds for their husband, Christ, being betrothed to him by sincere faith.

These follow the Lord, wherever he goes. That is to say, they care for no person but Christ; they desire no person but Christ. In him they place all their help, all their comfort, all their joy, all their salvation; to him alone they always give respect, in him they know themselves to be complete, which one and alone is to them all things. Moreover, wherever Christ calls the faithful by doctrine and example, even to death and most cruel slaughter, they follow willingly and cheerfully. By this means, they can never be separated from him in the world to come. For wherever Christ is, there is also Christ’s minister, as he himself witnessed in the twelfth and fourteenth chapters of John.

\index{Apocalypse (Book of)}They are also redeemed from humankind, truly delivered through the grace of Christ, that they should not follow this corrupt and unclean world, by all manner of pollution. For Christ by his Spirit and word calls his out of this world, that although in body we are conversant in the world, yet we should with all our mind abhor the world and the things that are therein. Furthermore, for this intent he has chosen and redeemed his from the bondage of humankind or of the world, that they should be first fruits to God the Father and to his Son. Which place the most godly and excellent learned man Dr. Francis Lambert expounding in his commentaries upon the Apocalypse: it is manifest, he says, by Leviticus 23:15, Numbers 18:15, and Deuteronomy 18:18, what first fruits are, and that they were gathered for the Lord and went to the high priest. But Christ is that high priest, to whom the spiritual first fruits appertain, namely, the godly and sanctified to God. These things are confirmed by the Apostle, who said that Christ gave himself for us, to the end he might redeem us from all iniquity and might purify us to himself, a people for good works. Therefore, the true faithful singularly apply this to godliness, that they may be the first fruits and a most excellent offering to the Lord, since they know themselves to be redeemed for this end, that all the rest of the time of their life they might serve God.

In their speech, there is no deceit. He does not say that no desire or evil impulse is found in the hearts of the faithful, but that there is no deceit in their speech. For although the faithful are troubled and vexed by the affections of the flesh, they love the truth so much that, to their knowledge, they will deceive no one. And they chiefly dissemble nothing that pertains to the confession of truth, nor use any deceit in the doctrine of the Gospel.

\index{Augustine, Saint}Moreover, they stand without blemish before the throne of God, not by their own merit, but by the sanctification of Christ, as Paul also affirms in chapter 5 to the Ephesians. He spoke appropriately of them standing before the throne. For Augustine said that our sanctification will ultimately be perfected in the world to come.

These are what I say are the true marks of the true faithful, and of the true church of Christ. Let every man search the secret corners of his heart and consider diligently in his mind whether he is marked with these signs. Let him earnestly pray to God that, if he feels them, the Lord would confirm them; if he does not feel them, that the Lord would impress them deeply in their minds.
